\section{CCD 境界スキームの係数導出}
\label{app:ccd_bc_derivation}
% ═══════════════════════════════════════════════════════════════════════════════

本付録は §\ref{sec:ccd_bc} の左境界スキームの係数を Taylor 展開により導出する．

\subsection{Equation-I 境界スキームの導出}
\label{app:ccd_bc_derivation_I}

境界点 $i=0$ において，片側スキーム（式 \eqref{eq:bc_general}）の未知数
$(\alpha, \beta, c_0, c_1, c_2, c_3)$ を決定する．

\paragraph{Step 1：$\alpha, \beta$ の決定}

$\fp{1}, \fpp{1}$ を $i=0$ 周りで展開して LHS に代入すると，
$\fpp{0}$ の係数消去条件 $\alpha + \beta = 0$（すなわち $\beta = -\alpha$）が得られる．

\paragraph{Step 2：$\Ord{h^5}$ 精度から $\alpha$ の一意決定}

$\beta = -\alpha$ 代入後，$S_k \equiv c_1 + 2^k c_2 + 3^k c_3$ として
5つの整合条件（定数項から $f^{(4)}_0$ まで）を立てる：
$c_0{+}c_1{+}c_2{+}c_3=0$，$S_1=1{+}\alpha$，$S_2=0$，$S_3=-3\alpha$，$S_4=-8\alpha$．

差分 $S_2{-}S_1$，$S_3{-}S_2$，$S_4{-}S_3$ を順次取ると：
\[
  2c_2+6c_3 = -(1+\alpha),\quad
  4c_2+18c_3 = -3\alpha,\quad
  8c_2+54c_3 = -5\alpha
\]
第1式の2倍を第2式から引くと $c_3 = (2-\alpha)/6$．
第1式に戻して $c_2 = -3/2$（$\alpha$ に依らず一定）．
第3式に代入：
\[
  -12 + 9(2-\alpha) = -5\alpha
  \quad\Longrightarrow\quad
  6 = 4\alpha
  \quad\Longrightarrow\quad
  \boxed{\alpha = \tfrac{3}{2}}
\]
$\alpha \neq 3/2$ では $\Ord{h^4}$ にとどまる—— $\alpha = 3/2$ は唯一の $\Ord{h^5}$ 選択である．

\paragraph{Step 3：残余係数の計算}

$\alpha = 3/2$，$c_2 = -3/2$ を用いて：
$c_3 = 1/12$，$c_1 = 21/4$，$c_0 = -23/6$．\qed

\subsection{Equation-II 境界スキームの導出}
\label{app:ccd_bc_derivation_II}

境界点 $i=0$ において，式 \eqref{eq:bcII_general} の未知数
$(\gamma, \delta, d_0, d_1, d_2, d_3)$ を決定する．

\paragraph{Step 1：$\gamma, \delta$ の決定}

$\fp{1}, \fpp{1}$ を $i=0$ 周りで展開して LHS に代入すると：
\[
  \text{LHS} = \gamma h\,\fp{0} + (1+\gamma+\delta)\fpp{0} + \cdots
\]
$\fp{0}$ の係数をゼロにする：$\gamma = 0$．
$\fpp{0}$ の係数を 1 に正規化：$1+\gamma+\delta = 1$，すなわち $\delta = 0$．
したがって LHS $= \fpp{0}$（片側関数値スキーム）．

\paragraph{Step 2：RHS の係数決定}

$f_k$（$k=0,1,2,3$）を $f_0$ 周りで展開し，$f_0/h^2$，$\fp{0}/h$，$\fpp{0}$，$f^{(3)}_0/h$ の
各係数整合条件を解く：
\begin{itemize}
  \item $\fpp{0}$係数$-$$\fp{0}/h$係数: $d_2+3d_3=1$
  \item $f^{(3)}_0/h$係数$-$$\fpp{0}$係数: $2d_2+9d_3=-1$
  \item $3d_3 = -3 \Rightarrow d_3=-1$，$d_2=4$，$d_1=-5$，$d_0=2$
\end{itemize}
結果：式 \eqref{eq:bcII_left} の $\fpp{0} = (2f_0-5f_1+4f_2-f_3)/h^2$（$\Ord{h^2}$ 精度）．\qed

\subsection{Equation-II 境界スキームの4次精度昇格（\texorpdfstring{$\Ord{h^4}$}{O(h\textasciicircum 4)}）}
\label{app:ccd_bc_derivation_II_h4}

4点公式（$\Ord{h^2}$）では全体の CCD d2 収束が境界で制限されるため，6点公式
（$\Ord{h^4}$）に昇格する（$n_\mathrm{pts} \ge 6$ のとき有効）．

\paragraph{導出}

$f_k$（$k=0,\ldots,5$）の Taylor 展開を $f_0$ 周りで行い，係数
$d_0,\ldots,d_5$ を連立方程式
\[
  \sum_{k=0}^{5} d_k \, f_k = \fpp{0} + \Ord{h^p}, \qquad p \ge 4
\]
の条件 $k=0,1,2,3,4$ から決定する（$k=5$ は $\Ord{h^4}$ 打ち切り誤差の確認）：

\begin{align}
  k=0:\quad & \sum d_k = 0         \label{eq:bcII_h4_k0}\\
  k=1:\quad & \sum k\,d_k = 0      \label{eq:bcII_h4_k1}\\
  k=2:\quad & \sum k^2 d_k = 2     \label{eq:bcII_h4_k2}\\
  k=3:\quad & \sum k^3 d_k = 0     \label{eq:bcII_h4_k3}\\
  k=4:\quad & \sum k^4 d_k = 0     \label{eq:bcII_h4_k4}\\
  k=5:\quad & \sum k^5 d_k \ne 0   \quad (\text{$\Ord{h^4}$ 確認})
\end{align}

唯一解：
\begin{equation}
  \fpp{0}
  = \frac{45f_0 - 154f_1 + 214f_2 - 156f_3 + 61f_4 - 10f_5}{12h^2}
  + \Ord{h^4}.
  \label{eq:bcII_left_h4}
\end{equation}

\paragraph{数値検証}

係数ベクトル $\bm{d} = [45, -154, 214, -156, 61, -10]/12$ に対して，
$\sum_{k=0}^{5} d_k k^p$ が $p=0,\ldots,4$ ですべて $0$（$p=2$ は $2$）を与え，
$p=5$ で $-548/12 \ne 0$ となることを確認した（機械精度）．

\paragraph{実装}

\texttt{ccd/ccd\_solver.py} の \texttt{\_boundary\_coeffs\_left(h, n\_pts)} において，
$n_\mathrm{pts} \ge 6$ のとき式~\eqref{eq:bcII_left_h4} を，$n_\mathrm{pts} < 6$ のとき
旧来の $\Ord{h^2}$ 公式をフォールバックとして使用する．右境界では
$f_0 \to f_N,\, f_1 \to f_{N-1},\ldots$ の鏡像公式を適用する．

\paragraph{収束効果}

昇格により d2 の全域 $L^\infty$ 収束が境界寄与によって制限されず，
内点スキームの $\Ord{h^6}$ 精度に近い収束率が得られる．
数値実験（\texttt{test\_ccd.py} Test~2）では平均スロープ ${\approx}4.8$ を確認した．\qed

